\documentclass[11pt,a4paper]{article}
\usepackage[utf8]{inputenc}
\usepackage[T1]{fontenc}
\usepackage[margin=1in]{geometry}
\usepackage{xcolor}
\usepackage{hyperref}
\definecolor{primarycolor}{RGB}{33,37,41}
\definecolor{secondarycolor}{RGB}{0,102,204}
\hypersetup{colorlinks=true,linkcolor=secondarycolor,urlcolor=secondarycolor}
\pagestyle{empty}
\setlength{\parindent}{0pt}
\setlength{\parskip}{0.9em}
\begin{document}

Dear Corvale Mobility Hiring Team,

I am applying for the Cloud Infrastructure / DevOps Contractor role. My background combines CI/CD ownership, repeatable test environments and a practical focus on making systems more reliable---experience I would bring to the infrastructure work supporting Corvale’s new train fleet.

At Odrava Health Software, I own GitLab CI pipelines for three product teams and built short-lived, per-merge-request environments deployed with Helm to a shared Kubernetes cluster. I reduced average pipeline time from 38 to 14 minutes and helped lower the flaky-test rate from 7\% to below 1\%. I also wrote the internal guide to how the CI system works, which is used in onboarding. These projects have involved improving workflows across teams and making infrastructure easier to use and understand.

Alongside my professional work, I built a personal three-node k3s homelab with GitOps deployment through Argo CD, monitoring with Prometheus, Grafana and Loki, and Terraform provisioning for a small Hetzner VPS. This gives me additional hands-on experience with infrastructure automation and observability. I am interested in Corvale’s challenge of shaping the platform that connects train systems with cloud services, where reliability, security and clear interfaces matter.

I would welcome the opportunity to discuss how my infrastructure-focused experience and continued development toward platform engineering could contribute to the fleet connectivity team. Thank you for your consideration.

Sincerely,\\
Marta Kowal-Brandt
\end{document}